\documentclass[10pt,a4paper]{amsart}
\usepackage{amsmath,amssymb,amsthm,graphicx}
\usepackage{verbatim}
\newenvironment{reference}{\comment}{\endcomment}
%\newenvironment{reference}{}{}
\newenvironment{slogan}{\comment}{\endcomment}
\newenvironment{history}{\comment}{\endcomment}

% For commutative diagrams we use Xy-pic
\usepackage[all]{xy}

% We use 2cell for 2-commutative diagrams.
\xyoption{2cell}
% Theorem environments.
%
\theoremstyle{plain}
\newtheorem{theorem}[subsection]{Theorem}
\newtheorem{proposition}[subsection]{Proposition}
\newtheorem{lemma}[subsection]{Lemma}

\theoremstyle{definition}
\newtheorem{definition}[subsection]{Definition}
\newtheorem{example}[subsection]{Example}
\newtheorem{exercise}[subsection]{Exercise}
\newtheorem{situation}[subsection]{Situation}

\theoremstyle{remark}
\newtheorem{remark}[subsection]{Remark}
\newtheorem{remarks}[subsection]{Remarks}

\numberwithin{equation}{subsection}

% Macros
%
\def\lim{\mathop{\mathrm{lim}}\nolimits}
\def\colim{\mathop{\mathrm{colim}}\nolimits}
\def\Spec{\mathop{\mathrm{Spec}}}
\def\Hom{\mathop{\mathrm{Hom}}\nolimits}
\def\Ext{\mathop{\mathrm{Ext}}\nolimits}
\def\SheafHom{\mathop{\mathcal{H}\!\mathit{om}}\nolimits}
\def\SheafExt{\mathop{\mathcal{E}\!\mathit{xt}}\nolimits}
\def\Sch{\mathit{Sch}}
\def\Mor{\mathop{\mathrm{Mor}}\nolimits}
\def\Ob{\mathop{\mathrm{Ob}}\nolimits}
\def\Sh{\mathop{\mathit{Sh}}\nolimits}
\def\NL{\mathop{N\!L}\nolimits}
\def\CH{\mathop{\mathrm{CH}}\nolimits}
\def\proetale{{pro\text{-}\acute{e}tale}}
\def\etale{{\acute{e}tale}}
\def\QCoh{\mathit{QCoh}}
\def\Ker{\mathop{\mathrm{Ker}}}
\def\Im{\mathop{\mathrm{Im}}}
\def\Coker{\mathop{\mathrm{Coker}}}
\def\Coim{\mathop{\mathrm{Coim}}}

% Boxtimes
%
\DeclareMathSymbol{\boxtimes}{\mathbin}{AMSa}{"02}

%
% Macros for moduli stacks/spaces
%
\def\QCohstack{\mathcal{QC}\!\mathit{oh}}
\def\Cohstack{\mathcal{C}\!\mathit{oh}}
\def\Spacesstack{\mathcal{S}\!\mathit{paces}}
\def\Quotfunctor{\mathrm{Quot}}
\def\Hilbfunctor{\mathrm{Hilb}}
\def\Curvesstack{\mathcal{C}\!\mathit{urves}}
\def\Polarizedstack{\mathcal{P}\!\mathit{olarized}}
\def\Complexesstack{\mathcal{C}\!\mathit{omplexes}}
% \Pic is the operator that assigns to X its picard group, usage \Pic(X)
% \Picardstack_{X/B} denotes the Picard stack of X over B
% \Picardfunctor_{X/B} denotes the Picard functor of X over B
\def\Pic{\mathop{\mathrm{Pic}}\nolimits}
\def\Picardstack{\mathcal{P}\!\mathit{ic}}
\def\Picardfunctor{\mathrm{Pic}}
\def\Deformationcategory{\mathcal{D}\!\mathit{ef}}
\usepackage{cleveref}
\crefname{originallemma}{original lemma}{original lemmas}
\Crefname{originallemma}{Original lemma}{Original lemmas}
\makeatletter
\newlabel{chow-lemma-factors}{{02TI}{source}{}{sourcepointer-1}{}}
\newlabel{chow-lemma-factors@cref}{{[originallemma][0][]02TI}{[1][0][]source}{}{}{}}
\makeatother
\crefname{originallemma}{original lemma}{original lemmas}
\Crefname{originallemma}{Original lemma}{Original lemmas}
\makeatletter
\newlabel{chow-lemma-proper-pushforward-rational-equivalence}{{02S2}{source}{}{sourcepointer-2}{}}
\newlabel{chow-lemma-proper-pushforward-rational-equivalence@cref}{{[originallemma][0][]02S2}{[1][0][]source}{}{}{}}
\makeatother
\crefname{originaldefinition}{original definition}{original definitions}
\Crefname{originaldefinition}{Original definition}{Original definitions}
\makeatletter
\newlabel{divisors-definition-strict-transform}{{080D}{source}{}{sourcepointer-3}{}}
\newlabel{divisors-definition-strict-transform@cref}{{[originaldefinition][0][]080D}{[1][0][]source}{}{}{}}
\makeatother
\crefname{originallemma}{original lemma}{original lemmas}
\Crefname{originallemma}{Original lemma}{Original lemmas}
\makeatletter
\newlabel{divisors-lemma-blowing-up-projective}{{02NS}{source}{}{sourcepointer-4}{}}
\newlabel{divisors-lemma-blowing-up-projective@cref}{{[originallemma][0][]02NS}{[1][0][]source}{}{}{}}
\makeatother
\title{Hironaka's example of a proper morphism of smooth complex threefolds that is Zariski-locally projective but not projective}
\author{The Stacks Project Authors}
\date{stacks-project:examples.tex:local-snapshot-20261012; GNU FDL 1.2 or later}
\begin{document}
\maketitle
\noindent\textit{Editorial note.} All delivered mathematical text is the author's own, copied from the pinned Stacks Project chapter file examples.tex; no statement or proof step was rewritten, reordered or removed.\par\medskip
\noindent\textit{Editorial note.} The counted proof is the author's complete argument: Hironaka example; the companion lemma lemma-non-effective-descent-projective reuses this construction and its equivalence conclusion.\par\medskip
\noindent\textit{Editorial note.} The label lemma-non-descending-property-projective names the counted statement, and the proof environment immediately following it in the source is delivered verbatim as its proof.\par\medskip
\noindent\textit{Editorial note.} The counted statement is carried by the author's own lemma environment in the chapter source.\par\medskip
\noindent\textit{Editorial note.} Revision r002: regenerated from the same spec after the r001 assembly failed at compile time (the author's citation of Hironaka via the bare key H had no bibliography backend in r001; the stale r001 job directory is retained as evidence).\par\medskip
\noindent\textit{Editorial note.} Bibliography support: the counted proof cites the letter-key H (the Stacks Project's global bibliography entry for Hartshorne, Algebraic Geometry, Graduate Texts in Mathematics 52, Springer, 1977, whose Appendix B Example B.3.4.1/B.3.4.2 are the Hironaka examples being imitated). A matching hironaka.bib asset is attached so the author's own citation resolves; no statement or proof text was altered.\par\medskip
\noindent\textit{Allowed standard prerequisites (editorial).}\begin{itemize}
\item Divisors, Definition divisors-definition-strict-transform: the strict transform of a closed subscheme under a blowup
\item Divisors, Lemma divisors-lemma-blowing-up-projective: a blowup morphism is projective
\item Chow Homology, Lemma chow-lemma-proper-pushforward-rational-equivalence and Lemma chow-lemma-factors: the degree of a line bundle on a proper scheme over a field is well defined on one-cycles modulo rational equivalence
\item The conics C given by x y equals z squared and w equals zero, and D given by x y equals w squared and z equals zero, in projective three-space over the complex numbers are smooth, meet transversally in the two reduced points P equals the point with coordinates one zero zero zero and Q equals the point with coordinates zero one zero zero, and each is isomorphic to the projective line
\end{itemize}\medskip
\ifdefined\hypertarget\hypertarget{sourcepointer-1}{}\fi
\noindent\textit{Source reference.} Original lemma 02TI: sources/stacks-chow/tags.txt, source lines 3697--3697. Prerequisite pointer invoked at 5434; it names a standard fact of a sibling chapter that the counted argument uses, and every step of the counted proof built on it is retained verbatim in the delivered body.\par\medskip
\ifdefined\hypertarget\hypertarget{sourcepointer-2}{}\fi
\noindent\textit{Source reference.} Original lemma 02S2: sources/stacks-chow/tags.txt, source lines 3645--3645. Prerequisite pointer invoked at 5433; it names a standard fact of a sibling chapter that the counted argument uses, and every step of the counted proof built on it is retained verbatim in the delivered body.\par\medskip
\ifdefined\hypertarget\hypertarget{sourcepointer-3}{}\fi
\noindent\textit{Source reference.} Original definition 080D: sources/stacks-chow/tags.txt, source lines 10178--10178. Prerequisite pointer invoked at 5398; it names a standard fact of a sibling chapter that the counted argument uses, and every step of the counted proof built on it is retained verbatim in the delivered body.\par\medskip
\ifdefined\hypertarget\hypertarget{sourcepointer-4}{}\fi
\noindent\textit{Source reference.} Original lemma 02NS: sources/stacks-chow/tags.txt, source lines 3491--3491. Prerequisite pointer invoked at 5405; it names a standard fact of a sibling chapter that the counted argument uses, and every step of the counted proof built on it is retained verbatim in the delivered body.\par\medskip
\begin{lemma}
\label{lemma-non-descending-property-projective}
The properties
\begin{enumerate}
\item[] $\mathcal{P}(f) =$``$f$ is projective'', and
\item[] $\mathcal{P}(f) =$``$f$ is quasi-projective''
\end{enumerate}
are not Zariski local on the base. A fortiori, they are not fpqc local
on the base.
\end{lemma}

\begin{proof}
Following Hironaka \cite[Example B.3.4.1]{H},
we define a proper morphism of smooth complex 3-folds
$f:V_Y\to Y$
which is Zariski-locally projective, but not projective. Since $f$ is proper
and not projective, it is also not quasi-projective.

\medskip\noindent
Let $Y$ be projective 3-space over the complex numbers ${\mathbf C}$.
Let $C$ and $D$ be smooth conics in $Y$ such
that the closed subscheme $C\cap D$ is reduced and consists
of two complex points $P$ and $Q$. (For example,
let $C=\{ [x,y,z,w]: xy=z^2, w=0\}$, $D=\{ [x,y,z,w]:
xy=w^2, z=0\}$, $P=[1,0,0,0]$,
and $Q=[0,1,0,0]$.) 
On $Y-Q$, first blow up the curve $C$, and then blow
up the strict transform of the curve $D$ (Divisors, Definition 
\ref{divisors-definition-strict-transform}). On $Y-P$, first blow up
the curve $D$, and then blow up the strict transform of the curve
$C$. Over $Y-P-Q$, the two varieties we have constructed are canonically
isomorphic, and so we can glue them over $Y-P-Q$. The result
is a smooth proper 3-fold $V_Y$ over ${\mathbf C}$. The morphism
$f:V_Y\to Y$ is proper and Zariski-locally projective (since
it is a blowup over $Y-P$ and over $Y-Q$), by Divisors,
Lemma \ref{divisors-lemma-blowing-up-projective}. We will show that
$V_Y$ is not projective over ${\mathbf C}$. That will imply that
$f$ is not projective.

\medskip\noindent
To do this, let $L$ be the inverse image in $V_Y$ of a complex point
of $C-P-Q$, and $M$ the inverse image of a complex point of $D-P-Q$.
Then $L$ and $M$ are isomorphic to the projective line
${\mathbf P}^1_{{\mathbf C}}$. 
Next, let $E$ be the inverse image in $V_Y$ of $C\cup D\subset Y$ in $V_Y$;
thus $E\rightarrow C\cup D$ is a proper morphism, with fibers
isomorphic to ${\mathbf P}^1$ over $(C\cup D)-\{P,Q\}$.
The inverse
image of $P$ in $E$ is a union of two lines $L_0$ and $M_0$, and we have
rational equivalences of cycles $L\sim L_0+M_0$ and $M\sim M_0$ on $E$
(using that $C$ and $D$ are isomorphic to ${\mathbf P}^1$).
Note the asymmetry resulting from the order in which we blew
up the two curves. Near $Q$, the opposite happens. So the inverse image
of $Q$ is the union of two lines $L_0'$ and $M_0'$, and we have
rational equivalences $L\sim L_0'$ and $M\sim L_0'+M_0'$ on $E$.
Combining these equivalences, we find that $L_0+M_0'\sim 0$
on $E$ and hence on $V_Y$. If $V_Y$ were projective over ${\mathbf C}$,
it would have
an ample line bundle $H$, which would have degree $> 0$ on all curves
in $V_Y$. In particular $H$ would have positive degree on $L_0+M_0'$,
contradicting that the degree of a line bundle is well-defined
on 1-cycles modulo rational equivalence on a proper scheme
over a field (Chow Homology,
Lemma \ref{chow-lemma-proper-pushforward-rational-equivalence}
and Lemma \ref{chow-lemma-factors}).
So $V_Y$ is not projective over ${\mathbf C}$.
\end{proof}

\bibliographystyle{plain}
\bibliography{hironaka}
\end{document}
